\documentclass{abntex2}
\usepackage[utf8]{inputenc}
\usepackage[T1]{fontenc}
\usepackage{geometry}
\usepackage{graphicx}
\usepackage{float}
\usepackage{hyperref}
\usepackage[english,indonesian]{babel}

\title{Optimasi Pemrosesan Template LaTeX Menggunakan Pendekatan Berbasis AI}
\author{Academic Support Agent}
\date{\today}

\begin{document}
\maketitle

\begin{abstract}

\end{abstract}

\chapter{PENDAHULUAN}


\chapter{TINJAUAN PUSTAKA}


% Original content (commented out):
% \documentclass[12pt,openright,oneside,a4paper,brazil]{abntex2}

% openany
% twoside - frente e verso
% oneside - somente frente

% ---
% Créditos:
% Leonardo Pedrosa - versão 21.05
% Tomáz Malta Conceição - versão 22.03
% ---

% ---
% Possibilita criação de Quadros e Lista de quadros.
% Ver https://github.com/abntex/abntex2/issues/176
%
\newcommand{\quadroname}{Quadro}
\newcommand{\listofquadrosname}{Lista de quadros}

\newfloat[chapter]{quadro}{loq}{\quadroname}
\newlistof{listofquadros}{loq}{\listofquadrosname}
\newlistentry{quadro}{loq}{0}

% configurações para atender às regras da ABNT
\setfloatadjustment{quadro}{\centering}
\counterwithout{quadro}{chapter}
\renewcommand{\cftquadroname}{\quadroname\space} 
\renewcommand*{\cftquadroaftersnum}{\hfill--\hfill}

\setfloatlocations{quadro}{hbtp} % Ver https://github.com/abntex/abntex2/issues/176
% ---



% \uppercase is not supported when a string is expanded for bookmarks. Either ignore the warning or add in the preamble:
\pdfstringdefDisableCommands{\let\uppercase\relax}

% Observações: para textos muito longos no arquivo abntex2.bib nesse exemplo aconteceu com a citação do SEBRAE. obs no arquivo dicas.tex
%\usepackage{etoolbox}
%\apptocmd{\thebibliography}{\raggedright}{}{}
% \apptocmd{\thebibliography}{\raggedjustifying}{}{}
% \apptocmd{\thebibliography}{\justifying}{}{}
\usepackage{etoolbox}
\apptocmd{\sloppy}{\hbadness 10000\relax}{}{}


%\usepackage[num]{abntex2cite}
% incluir para utilizar o abntex2 na citação com números

%\usepackage[num,bibjustif]{abntex2cite}
\usepackage[alf,bibjustif]{abntex2cite}
% incluir para utilizar o abntex2 na citação
% bibjustif - deixa as referências justificadas
% alf - para referenciar com nomes e anos, a ordem no fianal é alfabética
% num - para colocar número nas referências e referenciar no final do texto na ordem que aparece no texto

\usepackage{verbatim}

\usepackage[brazil]{babel}
\usepackage[utf8]{inputenc}

\usepackage[T1]{fontenc}
% se for utilizar o comando \hyphenation com palavras acentuadas é preciso retirar o comentário do comando acima, porém a qualidade da escrita quando gera-se um pdf parece que fica um pouco prejudicada. Para a qualidade não ficar prejudicada utiliza-se o comando a seguir
\usepackage{ae}

% usado por: abntex2-fichacatalografica.tex, para colocar o número de páginas
\usepackage{lastpage}

\usepackage{pdfpages}
% para inserir páginas em pdf
% utilizado para inserir a folha de aprovação assinada

\usepackage{float}
% para posicionar a figura no local que tem interesse utilizando o "H"
% ou seja: \begin{figure}[H]

% Você pode customizar a fonte dos títulos dos capítulos alterando os comandos como no exemplo a seguir, para que seja utilizada a fonte Computer Modern com tamanho maior do que o utilizado por padrão:
% para que título e subtítulo fiquem com a mesma letra
\renewcommand{\ABNTEXchapterfont}{\fontfamily{cmr}\fontseries{b}\selectfont}
\renewcommand{\ABNTEXchapterfontsize}{\Large}
% \Large para diminuir o tamanho dos títulos dos capítulos que estavam com \HUGE

% altera o tamanho das seções, subseções e capítulos
\renewcommand{\ABNTEXsectionfontsize}{\normalsize}
\renewcommand{\ABNTEXsubsectionfontsize}{\normalsize}

% altera o estilo do texto dos capítulos
% the memoir class página 385
% \chapterstyle{article}

% Indenta o primeiro parágrafo de cada seção.
\usepackage{indentfirst}

% backref permite que a bibliografia indique quantas vezes e em quais páginas a citação ocorreu
% \usepackage[brazilian]{backref}
% \usepackage{backref}
% \usepackage[backref=page]{hyperref}
% \renewcommand{\backrefpagesname}{Citado na(s) página(s):~}

%\usepackage{titletoc}
%\dottedcontents{section}[0em]{\bfseries}{2.9em}{1pc}
%\dottedcontents{subsection}[0em]{}{3.3em}{1pc}
%\renewcommand{\contentsname}{\centering Contents}



%Hyphenations, caso seja necessário acrescentar alguma...
\hyphenation{a-de-que-do}
\hyphenation{cons-tru-ção}
\hyphenation{cor-res-pon-dên-cia}
\hyphenation{contribuição}
\hyphenation{SILVA}
\hyphenation{NOSELLA}

\titulo{Inserir o título aqui...}
\autor{Inserir o nome do autor aqui acentuação......}
\data{Inserir o número do ano aqui...}
\instituicao{Universidade Federal de Ouro Preto}
\local{Ouro Preto}
\preambulo{Monografia apresentada ao Curso de Engenharia de Produção da Universidade Federal de Ouro Preto como parte dos requisitos para a obtenção do Grau de Engenheiro de Produção.}
\tipotrabalho{Monografia}
\orientador{Cristiano Luís Turbino de França e Silva}
\coorientador{Tiago}

\begin{document}

% elementos pré-textuais
\begin{frontmatter}
    \input{tex/frontmatter}
\end{frontmatter}

% Inicio do conteúdo do corpo do trabalho

%\begin{textual}
\begin{mainmatter}
    \input{tex/introducao}
    \input{tex/referencial_teorico}
    \input{tex/metodologia}
    \input{tex/apresentacao_resultados}
    \input{tex/consideracoes_finais}
    \input{tex/dicas}
\end{mainmatter}
%\end{textual}

\bibliography{abntex2}

% ---
% Inicia os anexos
\begin{anexosenv}
% Imprime uma página indicando o início dos anexos
\partanexos
    \input{tex/anexo}
\end{anexosenv}
% Termina os anexos
% ---

% incluir para indicar a formatação das referências, no final do trabalho
% \bibliographystyle{abntex2-alf}

% \bibliographystyle{abnt-num}
% Dessa forma, as citações no texto seguirão o padrão do LATEX, porém a bibliografia será formatada conforme o padrão brasileiro.

\end{document}

\end{document}